柱坐标和球坐标的解要先满足拉普拉斯方程,再套边界.勒让德展开的系数要用正交关系算出来.

\begin{problemset}
\renewcommand{\theenumi}{\thechapter.\arabic{enumi}}
\renewcommand{\labelenumi}{\theenumi.}
\item 由位置矢量算出柱坐标和球坐标的尺度因子.说明 $\phi$ 方向的弧长在柱坐标里是 $\rho\,\mathrm{d}\phi$,在球坐标里是 $r\sin\theta\,\mathrm{d}\phi$.

\item 均匀电场 $\mathbf{E}_0$ 垂直于无限长导体圆柱的轴,柱面接地.在柱外求电位,并计算柱面上的电场.说明它为什么是 $2E_0\cos\phi$.

\item 圆盘 $\rho<a$ 的边界电位是 $V_0\cos\phi$.求盘内电位,并在 $\rho=a$ 上核对边界值.

\item 有限圆柱 $0\leq\rho<a$, $0\leq z<b$,侧面和底面电位为零,顶面 $z=b$ 的电位是 $f(\rho)$.写出用 $J_0$ 表示的级数,并把系数写成积分.

\item 用罗德里格斯公式算出 $P_0,P_1,P_2$.验证它们满足勒让德方程,并且 $P_\ell(1)=1$.

\item 均匀电场中的接地导体球,写出球外电位和球面电场.说明表面场是 $3E_0\cos\theta$.

\item 球面 $r=a$ 上的电位是 $V_0\cos^2\theta$.把它展成 $P_0$ 和 $P_2$,分别写出球内和球外的电位.

\item 用泊松核说明圆盘中心的电位是边界值的圆周平均.再对 $g(\theta)=\cos\theta$ 算出盘内电位,与分离变量的结果比较.

\item 圆盘诺伊曼问题 $u_r(a,\theta)=h(\theta)$.证明 $h$ 的平均值必须是零,并说明解为什么可以加任意常数.

\item 扇形 $r<3$, $|\theta|<\pi/3$,两条边上 $u=0$,圆弧 $r=3$ 上 $u=g(\theta)$.写出分离变量后的级数,径向用哪个幂次.

\item 半环 $1<r<2$, $0<\theta<\pi$,三条边上 $u=0$, $r=2$ 上 $u=\pi-\theta$.把系数积分出来.

\item 圆膜 $r<a$ 固定在边界上,做小振动.说明频率由 $J_m(ka)=0$ 决定,并指出 $m\geq 1$ 时每个频率至少有两个独立的角向模式.

\item 写出 $\ell=0,1,2$ 的一组齐次调和多项式,并直接验证 $\ell=2$ 的五个多项式调和.说明 $2\ell+1$ 从何而来.
\end{problemset}
